% =============================================================================
%                    WEEK 4: COOPERATIVE GAMES (PART I)
% =============================================================================
\section{Week 4: Cooperative Games (Part I)}

% -----------------------------------------------------------------------------
%                              4.1 TOPIC SUMMARY
% -----------------------------------------------------------------------------
\subsection{Topic Summary}

\begin{keybox}[Learning Objectives]
By the end of this week, you will be able to:
\begin{enumerate}
  \item Explain the fundamental concepts of cooperative game theory, including \textbf{coalitions}, \textbf{transferable utility}, and the difference from non-cooperative settings.
  \item Model real-world scenarios as cooperative games, identifying possible coalitions and collective payoff structures.
  \item Analyse solution concepts such as the \textbf{core}, \textbf{Shapley value}, and \textbf{Banzhaf index}, including their properties and interpretation.
  \item Determine when stable and fair coalition structures exist, and evaluate conditions for allocations to be in the core.
  \item Calculate the Shapley value and Banzhaf index using mathematical reasoning.
  \item Assess the fairness, efficiency, and stability of division outcomes in cooperative environments.
\end{enumerate}
\end{keybox}

\separator

\subsubsection*{The Big Picture}

This week shifts from non-cooperative settings (where players act individually) to \textbf{cooperative game theory}, where players can form \textbf{binding agreements} and work in \textbf{coalitions}. The key premise: cooperation is not altruism---it emerges because players find it mutually beneficial to commit to agreements.

Two central questions drive the theory:
\begin{itemize}
  \item \textbf{Stability:} Which outcomes are immune to coalitional deviations? $\to$ \textbf{The Core}
  \item \textbf{Fairness:} How should payoffs reflect each player's contribution? $\to$ \textbf{Shapley Value}, \textbf{Banzhaf Index}
\end{itemize}

\separator

\subsubsection*{What Examiners Like to Ask}

\begin{itemize}
  \item Write the \textbf{formal definition} of a TU game $(N, v)$ and explain the characteristic function.
  \item Given a characteristic function, determine whether the game is \textbf{superadditive}, \textbf{convex}, or \textbf{simple}.
  \item Compute the \textbf{core} of a game by writing down the constraint system and finding feasible payoff vectors.
  \item Show that a given game has an \textbf{empty core} (e.g., 3-player majority game).
  \item Prove that \textbf{convex games} have non-empty cores.
  \item Prove that a \textbf{simple game} has a non-empty core iff it has a veto player.
  \item Compute the \textbf{Shapley value} for all players (enumerate permutations, compute marginal contributions).
  \item Compute the \textbf{Banzhaf index} and compare with the Shapley value.
  \item State and apply the \textbf{four axiomatic properties} of the Shapley value.
  \item Identify \textbf{null players} and \textbf{veto players} in simple/weighted voting games.
\end{itemize}

\separator

% -----------------------------------------------------------------------------
%           4.2 OVERVIEW + CORE CONCEPTS + EXTENDED THEORY
% -----------------------------------------------------------------------------
\subsection{Overview + Core Concepts + Extended Theory}

\subsubsection*{Roadmap}
\begin{enumerate}
  \item Cooperative vs.\ non-cooperative games; binding agreements
  \item TU vs.\ NTU games; characteristic vs.\ partition function games
  \item Formal definition of TU games: $(N, v)$
  \item Outcomes, coalition structures, and imputations
  \item Special classes: simple, superadditive, convex games
  \item The Core: definition, examples, emptiness
  \item Core of convex games and simple games
  \item Relaxations: $\varepsilon$-core, least core, cost of stability
  \item Shapley value: definition, computation, properties
  \item Banzhaf index
\end{enumerate}

\bigseparator

% --- Cooperative Games Introduction ---
\subsubsection{From Non-Cooperative to Cooperative Games}

\begin{defbox}[Cooperative (Coalitional) Game]
In \textbf{cooperative games}:
\begin{itemize}
  \item Agents can benefit by cooperating
  \item \textbf{Binding agreements} are possible (unlike Prisoner's Dilemma)
  \item Actions are taken by \textbf{coalitions} of agents
  \item Focus shifts from individual strategies to \textit{which coalitions form} and \textit{how payoffs are divided}
\end{itemize}
\end{defbox}

\begin{defbox}[TU vs.\ NTU Games]
\begin{center}
\renewcommand{\arraystretch}{1.3}
\begin{tabular}{@{} l p{5.5cm} p{5.5cm} @{}}
\toprule
 & \textbf{TU (Transferable Utility)} & \textbf{NTU (Non-Transferable Utility)} \\
\midrule
Payoff & Given to the group, then divided & Directly to individual members \\
Side payments & Feasible (money, goods) & Not possible \\
Example & Farmers splitting profit & Researchers getting promotions \\
Key question & How to split the surplus? & Which coalitions form? \\
\bottomrule
\end{tabular}
\end{center}
\end{defbox}

\begin{defbox}[CFG vs.\ PFG]
\begin{itemize}
  \item \textbf{Characteristic Function Game (CFG):} A coalition's payoff depends \textit{only} on its own composition. Each coalition is identified with a single number $v(C)$.
  \item \textbf{Partition Function Game (PFG):} A coalition's payoff depends on the \textit{entire partition}---what other coalitions are doing matters (e.g., market competition).
\end{itemize}

Most classical cooperative game theory (core, Shapley value) was developed for CFGs.
\end{defbox}

\separator

% --- TU Games Formalised ---
\subsubsection{Transferable Utility Games: Formal Definition}

\begin{defbox}[TU Game]
A \textbf{transferable utility game} is a pair $(N, v)$ where:
\begin{itemize}
  \item $N = \{1, \ldots, n\}$ is a set of \textbf{players}
  \item $v: 2^N \to \mathbb{R}$ is the \textbf{characteristic function}
  \item For each coalition $C \subseteq N$, $v(C)$ is the total payoff the members of $C$ can earn together
\end{itemize}

\textbf{Standard assumptions:}
\begin{itemize}
  \item \textbf{Normalised:} $v(\emptyset) = 0$
  \item \textbf{Non-negative:} $v(C) \geq 0$ for all $C \subseteq N$
  \item \textbf{Monotone:} $C \subseteq D \implies v(C) \leq v(D)$
\end{itemize}

$N$ itself is called the \textbf{grand coalition}.
\end{defbox}

\begin{defbox}[Outcome of a TU Game]
An \textbf{outcome} of $(N, v)$ is a pair $(\text{CS}, \mathbf{x})$ where:
\begin{itemize}
  \item $\text{CS} = (C_1, \ldots, C_k)$ is a \textbf{coalition structure}---a partition of $N$
  \item $\mathbf{x} = (x_1, \ldots, x_n)$ is a \textbf{payoff vector} satisfying:
  \begin{itemize}
    \item $x_i \geq 0$ for all $i \in N$
    \item $x(C) = \sum_{i \in C} x_i = v(C)$ for each $C$ in CS
  \end{itemize}
\end{itemize}

An outcome is an \textbf{imputation} if it also satisfies \textbf{individual rationality}: $x_i \geq v(\{i\})$ for all $i$.
\end{defbox}

\begin{exbox}[Running Example: Ice-Cream Game]
Three children: Charlie ({\pounds}4), Marcy ({\pounds}3), Pattie ({\pounds}2). Ice-cream tubs: 1kg for {\pounds}9, 750g for {\pounds}7, 500g for {\pounds}5. Children value ice-cream, not money.

\textbf{Characteristic function:}
\begin{center}
\renewcommand{\arraystretch}{1.2}
\begin{tabular}{@{} l c c c c c c c @{}}
\toprule
Coalition & $\emptyset$ & $\{C\}$ & $\{M\}$ & $\{P\}$ & $\{C,P\}$ & $\{M,P\}$ & $\{C,M\}$ \\
\midrule
$v(C)$ & 0 & 0 & 0 & 0 & 500 & 500 & 750 \\
\bottomrule
\end{tabular}
\end{center}
$v(\{C, M, P\}) = 1000$

\textbf{Some outcomes:}
\begin{itemize}
  \item $\text{CS} = \{\{C\}, \{M,P\}\}$; $x = (0, 150, 350)$
  \item $\text{CS} = \{N\}$; $x = (334, 333, 333)$
  \item $\text{CS} = \{N\}$; $x = (0, 800, 200)$ \quad (feasible but arguably unfair)
\end{itemize}
\end{exbox}

\separator

% --- Examples ---
\subsubsection{Important Game Types}

\begin{defbox}[Weighted Voting Games (WVG)]
Players have weights $w_1, \ldots, w_n$ and a quota $q$:
\[
v(C) = \begin{cases} 1 & \text{if } w(C) = \sum_{i \in C} w_i \geq q \quad (\text{winning}) \\ 0 & \text{otherwise} \quad (\text{losing}) \end{cases}
\]

\textbf{Notation:} $\text{WVG}[q;\; w_1, w_2, \ldots, w_n]$
\end{defbox}

\begin{exbox}[UK Elections as WVGs]
\begin{center}
\renewcommand{\arraystretch}{1.2}
\begin{tabular}{@{} l c c c @{}}
\toprule
\textbf{Year} & \textbf{Conservatives} & \textbf{Veto player?} & \textbf{Null players?} \\
\midrule
2010 & 306 seats & None & None \\
2019 & 365 seats & C (veto) & All others (null) \\
\bottomrule
\end{tabular}
\end{center}
(Quota = 326 in all cases, 650 total seats)
\end{exbox}

\begin{defbox}[Other Game Types]
\begin{itemize}
  \item \textbf{Network flow games:} Players are edges; $v(C) =$ max $s$--$t$ flow through edges in $C$
  \item \textbf{Graph-induced games:} Players are vertices; $v(C) =$ total weight of internal edges (can be negative)
\end{itemize}
\end{defbox}

\bigseparator

% --- Special Classes ---
\subsubsection{Special Classes of TU Games}

\begin{defbox}[Simple Games]
A game $(N, v)$ is \textbf{simple} if:
\begin{itemize}
  \item $v(C) \in \{0, 1\}$ for all $C \subseteq N$
  \item $v$ is monotone
\end{itemize}

Coalitions are \textbf{winning} ($v = 1$) or \textbf{losing} ($v = 0$).

\textbf{Convention:} In simple games, assume the grand coalition forms (outcome = payoff vector).
\end{defbox}

\begin{defbox}[Null and Veto Players (Simple Games)]
\begin{itemize}
  \item Player $i$ is a \textbf{null player} if: $v(C \cup \{i\}) = v(C)$ for all $C \subseteq N \setminus \{i\}$

  (Adding $i$ never changes any coalition's value)

  \item Player $i$ is a \textbf{veto player} if: $v(C) = 0$ for all $C \subseteq N \setminus \{i\}$

  (Equivalently, $v(N \setminus \{i\}) = 0$---without $i$, no coalition can win)
\end{itemize}
\end{defbox}

\separator

\begin{defbox}[Superadditive Games]
A game $(N, v)$ is \textbf{superadditive} if:
\[
v(S \cup T) \geq v(S) + v(T) \quad \text{for any disjoint } S, T \subseteq N
\]

\textbf{Interpretation:} Two coalitions can always merge without losing value $\Rightarrow$ grand coalition always forms.

\textbf{Convention:} In superadditive games, outcomes are identified with payoff vectors $\mathbf{x}$ satisfying $x(N) = v(N)$.
\end{defbox}

\begin{defbox}[Convex Games (Supermodularity)]
A function $f: 2^N \to \mathbb{R}$ is \textbf{supermodular} if $f(\emptyset) = 0$ and:
\[
f(A \cup B) + f(A \cap B) \geq f(A) + f(B) \quad \text{for all } A, B \subseteq N
\]

A game $(N, v)$ is \textbf{convex} if $v$ is supermodular.

\textbf{Key property (increasing marginal contributions):}
If $T \subset S$ and $i \notin S$, then:
\[
v(T \cup \{i\}) - v(T) \leq v(S \cup \{i\}) - v(S)
\]

A player is \textbf{more useful} when joining a \textbf{bigger} coalition.

\textbf{Hierarchy:} Convex $\Rightarrow$ Superadditive (but not vice versa).
\end{defbox}

\begin{exbox}[Properties of Game Types]
\begin{center}
\renewcommand{\arraystretch}{1.3}
\begin{tabular}{@{} l c c @{}}
\toprule
\textbf{Game Type} & \textbf{Superadditive?} & \textbf{Convex?} \\
\midrule
Ice-cream & Yes & No \\
Graph-induced & Iff all edge weights $\geq 0$ & Iff all edge weights $\geq 0$ \\
Network flow & Yes & Not necessarily \\
WVG (general) & Not always & Not always \\
WVG with $q > w(N)/2$ & Yes & Not necessarily \\
\bottomrule
\end{tabular}
\end{center}

\textbf{WVG counterexample for superadditivity:} $N = \{1,2\}$, $w_1 = w_2 = 1$, $q = 1$. Each singleton wins alone, so $v(\{1\}) + v(\{2\}) = 2 > 1 = v(N)$.

\textbf{WVG counterexample for convexity:} $N = \{1,2,3\}$, $w_i = 1$, $q = 2$. Then $v(\{1,2\}) + v(\{2,3\}) = 2 > 1 = v(\{1,2,3\}) + v(\{2\})$.
\end{exbox}

\bigseparator

% --- The Core ---
\subsubsection{The Core}

\begin{defbox}[The Core]
The \textbf{core} $\mathcal{C}(G)$ of a characteristic function game $G = (N, v)$ is the set of all outcomes $(\text{CS}, \mathbf{x})$ such that:
\[
x(C) \geq v(C) \quad \text{for every } C \subseteq N
\]

\textbf{Interpretation:} No coalition can deviate and earn more on its own---every coalition gets at least what it could guarantee itself.

For \textbf{superadditive} games, the core simplifies to all payoff vectors $\mathbf{x}$ with:
\begin{itemize}
  \item $x_i \geq 0$ for all $i$
  \item $x(N) = v(N)$
  \item $x(C) \geq v(C)$ for all $C \subseteq N$
\end{itemize}
\end{defbox}

\begin{thmbox}[Core Stability Implies Efficiency]
If $(\text{CS}, \mathbf{x})$ is in the core, then CS maximises social welfare:
\[
v(\text{CS}) \geq v(\text{CS}') \quad \text{for every coalition structure CS}' \text{ over } N
\]

\textbf{Proof sketch:} Suppose $v(\text{CS}) < v(\text{CS}')$. Then $\sum_{C' \in \text{CS}'} x(C') = v(\text{CS}) < v(\text{CS}') = \sum_{C' \in \text{CS}'} v(C')$. But core stability requires $x(C') \geq v(C')$ for each $C'$, giving $\sum x(C') \geq \sum v(C')$---contradiction.
\end{thmbox}

\begin{exbox}[Core of the Ice-Cream Game]
Since the game is superadditive, $\text{CS} = \{N\}$. Core constraints:
\begin{align*}
x_C \geq 0,\; x_M \geq 0,\; x_P &\geq 0 \\
x_C + x_M + x_P &= 1000 \\
x_C + x_M &\geq 750 \\
x_C + x_P &\geq 500 \\
x_M + x_P &\geq 500
\end{align*}

\textbf{Solutions include:} $(400, 400, 200)$, $(500, 400, 100)$, $(500, 500, 0)$.

Note: $(350, 350, 300)$ is \textbf{not} in the core because $x_C + x_M = 700 < 750$.
\end{exbox}

\begin{exbox}[Empty Core: 3-Player Majority Game]
$\text{WVG}[2;\; 1, 1, 1]$: $v(C) = 1$ if $|C| \geq 2$; $v(C) = 0$ if $|C| \leq 1$.

\textbf{The core is empty:}
\begin{itemize}
  \item If $\text{CS} = \{\{1\},\{2\},\{3\}\}$: grand coalition deviates (all get 0, but $v(N) = 1$)
  \item If $\text{CS} = \{\{1,2\},\{3\}\}$: one of $\{1,2\}$ gets $< 1$, can deviate with player 3
  \item If $\text{CS} = \{N\}$: some $x_i > 0$, so $x(N \setminus \{i\}) < 1 = v(N \setminus \{i\})$
\end{itemize}

No outcome satisfies all core constraints simultaneously.
\end{exbox}

\separator

% --- Core of Special Classes ---
\subsubsection{Core of Special Game Classes}

\begin{thmbox}[Core of Convex Games]
Every \textbf{convex game} has a non-empty core.

\textbf{Constructive proof:} Order players as $1, \ldots, n$. Set:
\[
x_i = v(\{1, \ldots, i\}) - v(\{1, \ldots, i-1\})
\]
(Pay each player their marginal contribution to the coalition of predecessors.)

Then $\sum x_i = v(N)$ (telescoping). For any $C = \{i, j, \ldots, s\}$ with $i < j < \ldots < s$:
\[
v(C) = \sum_{\text{members}} [\text{marginal to predecessors in } C] \leq \sum_{\text{members}} x_k
\]
by supermodularity (marginal contribution to a subset $\leq$ marginal contribution to a superset).

\textbf{Complexity:} A core outcome can be constructed in $O(n)$ queries to $v$.
\end{thmbox}

\begin{thmbox}[Core of Simple Games]
A simple game has a non-empty core \textbf{if and only if} it has a \textbf{veto player}.

\textbf{Proof ($\Leftarrow$):} If $i$ is a veto player, set $x_i = 1$, $x_k = 0$ for $k \neq i$. For any $C$: if $i \in C$, then $x(C) = 1 \geq v(C)$; if $i \notin C$, then $v(C) = 0$.

\textbf{Proof ($\Rightarrow$):} Suppose no veto player. Take any payoff vector $\mathbf{x}$ with $x(N) = 1$. Then some $x_i > 0$, so $x(N \setminus \{i\}) < 1$. But since $i$ is not a veto player, $v(N \setminus \{i\}) = 1$---core constraint violated.
\end{thmbox}

\begin{thmbox}[Corollary: Core Characterisation for Simple Games]
In a simple game, $\mathbf{x}$ is in the core iff $x_i = 0$ for every non-veto player $i$.

Checking core non-emptiness requires a single query: compute $v(N \setminus \{i\})$ for each $i$.
\end{thmbox}

\begin{tipbox}[Exam Tip: WVG with Coalition Structures]
If coalition structures (not just grand coalition) are allowed in WVGs, deciding core non-emptiness becomes \textbf{NP-hard} (reduction from PARTITION).

Example: $\text{WVG}[2;\; 1,1,1,1]$---the outcome $(\{\{1,2\},\{3,4\}\}, (\frac{1}{2}, \frac{1}{2}, \frac{1}{2}, \frac{1}{2}))$ is in the core, even though there is no veto player.
\end{tipbox}

\separator

% --- Core Relaxations ---
\subsubsection{Core Relaxations}

\begin{defbox}[$\varepsilon$-Core]
An outcome $\mathbf{x}$ is in the \textbf{$\varepsilon$-core} of a superadditive game $G$ (for $\varepsilon \geq 0$) if:
\[
x(C) \geq v(C) - \varepsilon \quad \text{for all } C \subseteq N
\]

Allows deviations only if the gain exceeds $\varepsilon$ (accounts for deviation costs).
\end{defbox}

\begin{defbox}[Least Core]
The \textbf{least core} of $G$ is the $\varepsilon^*$-core where:
\[
\varepsilon^*(G) = \inf\{\varepsilon \mid \varepsilon\text{-core of } G \text{ is non-empty}\}
\]

Can be computed via LP:
\[
\min \varepsilon \quad \text{s.t.} \quad x_i \geq 0,\; x(N) = v(N),\; x(C) \geq v(C) - \varepsilon \;\forall C \subseteq N
\]

The least core is \textbf{always non-empty}.
\end{defbox}

\begin{defbox}[Cost of Stability]
The \textbf{cost of stability} of $G = (N, v)$ is the smallest $\Delta \geq 0$ such that the game $G^\Delta = (N, v^\Delta)$ with $v^\Delta(N) = v(N) + \Delta$ and $v^\Delta(C) = v(C)$ for $C \subsetneq N$ has a non-empty core.

Computed via LP: $\min \Delta$ s.t.\ $x_i \geq 0$, $x(N) = v(N) + \Delta$, $x(C) \geq v(C)$ for all $C \subseteq N$.
\end{defbox}

\bigseparator

% --- Shapley Value ---
\subsubsection{Shapley Value}

\begin{defbox}[Shapley Value]
Given a game $G = (N, v)$ with $|N| = n$, the \textbf{Shapley value} of player $i$ is:
\[
\phi_i(G) = \frac{1}{n!} \sum_{\pi \in \Pi_N} \Delta^\pi_G(i)
\]

where $\Pi_N$ is the set of all permutations of $N$, and the \textbf{marginal contribution} of $i$ w.r.t.\ permutation $\pi$ is:
\[
\Delta^\pi_G(i) = v(S_\pi(i) \cup \{i\}) - v(S_\pi(i))
\]
with $S_\pi(i) = \{j \in N \mid \pi(j) < \pi(i)\}$ (predecessors of $i$ in $\pi$).

\textbf{Probabilistic interpretation:} $\phi_i$ is $i$'s expected marginal contribution when players arrive in a uniformly random order.
\end{defbox}

\begin{exbox}[Shapley Value: Ice-Cream Game (Pattie)]
Six permutations of $\{C, M, P\}$:

\begin{center}
\renewcommand{\arraystretch}{1.2}
\begin{tabular}{@{} l l r @{}}
\toprule
\textbf{Permutation} & \textbf{Predecessors of $P$} & $\Delta^\pi(P)$ \\
\midrule
$CMP$ & $\{C, M\}$ & $v(N) - v(\{C,M\}) = 250$ \\
$MCP$ & $\{M, C\}$ & $v(N) - v(\{C,M\}) = 250$ \\
$CPM$ & $\{C\}$ & $v(\{C,P\}) - v(\{C\}) = 500$ \\
$MPC$ & $\{M\}$ & $v(\{M,P\}) - v(\{M\}) = 500$ \\
$PCM$ & $\emptyset$ & $v(\{P\}) - v(\emptyset) = 0$ \\
$PMC$ & $\emptyset$ & $v(\{P\}) - v(\emptyset) = 0$ \\
\bottomrule
\end{tabular}
\end{center}

$\phi_P = \frac{250 + 250 + 500 + 500 + 0 + 0}{6} = 250$
\end{exbox}

\begin{exbox}[Shapley Value Using Symmetry]
Game: $v(N) = 1$, $v(C) = 0$ for all $C \subsetneq N$, with $|N| = n$.

\textbf{Method 1 (direct):} $\Delta^\pi(i) = 1$ only if $i$ is last in $\pi$. There are $(n-1)!$ such permutations.
$\phi_i = \frac{(n-1)!}{n!} = \frac{1}{n}$

\textbf{Method 2 (properties):} All players are symmetric $\Rightarrow$ equal Shapley values. By efficiency: $n \cdot \phi_i = v(N) = 1 \Rightarrow \phi_i = \frac{1}{n}$.
\end{exbox}

\separator

% --- Shapley Properties ---
\subsubsection{Shapley Value: Axiomatic Properties}

\begin{thmbox}[The Four Properties of the Shapley Value]
The Shapley value is the \textbf{unique} payoff division scheme satisfying:

\begin{center}
\renewcommand{\arraystretch}{1.4}
\begin{tabular}{@{} c l p{7.5cm} @{}}
\toprule
\textbf{\#} & \textbf{Property} & \textbf{Statement} \\
\midrule
1 & \textbf{Efficiency} & $\sum_{i \in N} \phi_i(G) = v(N)$ \\
2 & \textbf{Null (Dummy) player} & If $v(C \cup \{i\}) = v(C)$ for all $C$, then $\phi_i(G) = 0$ \\
3 & \textbf{Symmetry} & If $v(C \cup \{i\}) = v(C \cup \{j\})$ for all $C \subseteq N \setminus \{i,j\}$, then $\phi_i = \phi_j$ \\
4 & \textbf{Additivity} & $\phi_i(G_1 + G_2) = \phi_i(G_1) + \phi_i(G_2)$ \\
\bottomrule
\end{tabular}
\end{center}

where $(G_1 + G_2) = (N, v^+)$ with $v^+(C) = v^1(C) + v^2(C)$.
\end{thmbox}

\begin{tipbox}[Exam Tip: Using Properties to Compute Shapley Value]
The four properties can simplify computation:
\begin{enumerate}
  \item \textbf{Symmetry} reduces distinct computations (symmetric players get equal shares)
  \item \textbf{Null player} eliminates players with zero contribution
  \item \textbf{Efficiency} pins down the total
  \item \textbf{Additivity} lets you decompose complex games into simpler ones
\end{enumerate}
\end{tipbox}

\bigseparator

% --- Banzhaf Index ---
\subsubsection{Banzhaf Index}

\begin{defbox}[Banzhaf Index]
Given $G = (N, v)$ with $|N| = n$, the \textbf{Banzhaf index} of player $i$ is:
\[
\beta_i(G) = \frac{1}{2^{n-1}} \sum_{C \subseteq N \setminus \{i\}} [v(C \cup \{i\}) - v(C)]
\]

\textbf{Key difference from Shapley:} Averages over all \textbf{coalitions} (of which there are $2^{n-1}$ not containing $i$), rather than all permutations ($n!$).
\end{defbox}

\begin{defbox}[Banzhaf vs.\ Shapley: Property Comparison]
\begin{center}
\renewcommand{\arraystretch}{1.3}
\begin{tabular}{@{} l c c @{}}
\toprule
\textbf{Property} & \textbf{Shapley} & \textbf{Banzhaf} \\
\midrule
Efficiency & \checkmark & \texttimes \\
Null (Dummy) player & \checkmark & \checkmark \\
Symmetry & \checkmark & \checkmark \\
Additivity & \checkmark & \checkmark \\
\bottomrule
\end{tabular}
\end{center}

The Banzhaf index does \textbf{not} redistribute $v(N)$---the sum $\sum \beta_i$ may differ from $v(N)$.
\end{defbox}

\begin{exbox}[Banzhaf Index Example]
Game: $v(N) = 1$, $v(C) = 0$ for $C \subsetneq N$, $|N| = n$.

For player $i$: only $C = N \setminus \{i\}$ gives marginal contribution 1; all other $2^{n-1} - 1$ coalitions give 0.

$\beta_i = \frac{1}{2^{n-1}}$

$\sum_{i \in N} \beta_i = \frac{n}{2^{n-1}}$, which for $n > 2$ is $< 1 = v(N)$.
\end{exbox}

\begin{defbox}[Normalised Banzhaf Index]
To ensure efficiency:
\[
\nu_i(G) = \frac{\beta_i(G)}{\sum_{j \in N} \beta_j(G)}
\]

This fixes efficiency but \textbf{breaks additivity}.
\end{defbox}

\begin{tipbox}[Shapley-Shubik Power Index]
In \textbf{simple games} (where $v(C) \in \{0,1\}$), the Shapley value and Banzhaf index measure the \textbf{power} of a player---the probability that the player can influence the outcome. The Shapley value in this context is called the \textbf{Shapley-Shubik power index}.
\end{tipbox}

\bigseparator

% --- Summary Table ---
\subsubsection{Summary: Key Concepts}

\begin{center}
\renewcommand{\arraystretch}{1.4}
\begin{tabular}{@{} l p{9.5cm} @{}}
\toprule
\textbf{Concept} & \textbf{Key Fact} \\
\midrule
TU game $(N, v)$ & Characteristic function $v: 2^N \to \mathbb{R}$; normalised, non-negative, monotone \\
Coalition structure & Partition of $N$; payoff vector distributes each coalition's value \\
Superadditive & Merging never loses value; grand coalition always forms \\
Convex & Supermodular $v$; increasing marginal contributions; always has non-empty core \\
Simple game & $v(C) \in \{0,1\}$; non-empty core iff veto player exists \\
Core & Stable outcomes: no coalition can deviate; implies efficiency \\
$\varepsilon$-core / least core & Relaxation allowing bounded instability; always non-empty \\
Shapley value & Unique: efficient + null player + symmetry + additivity \\
Banzhaf index & Averages over coalitions (not permutations); not efficient \\
\bottomrule
\end{tabular}
\end{center}

\bigseparator
